\section{A Bernoulli cancellation law for every trigamma power}
\label{tp:sec:allpowers}

The quartic report asks for a cancellation mechanism at the fifth and
higher powers, a fixed arithmetic coordinate space at each power, and a
generator with its constants specified \cite[Research question R4]{RepoQuartic}.
We give one recurrence for all powers.  It explains the earlier quadratic,
cubic, and quartic cancellations and yields new fifth- and sixth-power
formulas.  The essential ingredient is a Bernoulli polynomial, rather than
a new summation theorem for each successive power.

\subsection{Regular moments as ordinary convergent sums}

Throughout this section let
\begin{equation}\label{tp:eq:tails}
 A_n=\psi'(n+1)=\zeta(2)-H_n^{(2)},\qquad x_n=n+\tfrac12,
 \qquad n\geq0.
\end{equation}
Bernoulli numbers have the convention $B_1=-1/2$.  For $d\geq1$ define
\begin{equation}\label{tp:eq:h}
 h_{d,j}=\sum_{j_1+\cdots+j_d=j}\prod_{\nu=1}^d B_{2j_\nu}(1/2),
 \qquad j_\nu\geq0.
\end{equation}
In particular $h_{d,0}=1$ and $h_{d,1}=-d/12$.  The centered
Hurwitz expansion \cite[\S25.11]{DLMF} gives
\begin{equation}\label{tp:eq:centeredasy}
 A_n^d\sim\sum_{j\geq0}h_{d,j}x_n^{-d-2j}.
\end{equation}
This is an asymptotic expansion; it is not a convergent Laurent series
at a fixed finite $n$.

\begin{proposition}[Ordinary realization of the regular parity]
\label{tp:prop:ordinary}
The function
\[
 F_d(s)=\sum_{n\geq0}A_n^d x_n^{-s},\qquad \Re s>1-d,
\]
has a meromorphic continuation to $\mathbb C$.  Its poles are exactly
$1-d-2j$, $j\geq0$, and their residues are $h_{d,j}$.  For every
$m\geq0$ with $m\equiv d\pmod2$, the value $M_{d,m}=F_d(-m)$ is regular
and is the absolutely convergent sum
\begin{equation}\label{tp:eq:ordinary}
 M_{d,m}=\sum_{n\geq0}\left(
 x_n^m A_n^d-
 \sum_{0\leq j\leq(m-d)/2}h_{d,j}x_n^{m-d-2j}\right).
\end{equation}
An inner sum with negative upper limit is empty.
\end{proposition}

\begin{proof}
Subtract a finite number of terms of \eqref{tp:eq:centeredasy} and add
the corresponding functions $h_{d,j}\zeta(s+d+2j,1/2)$.  Increasing
the subtraction order gives normally convergent remainders in
successively larger half-planes.  The meromorphic continuations agree
on overlaps, since they agree with the defining series in its initial
half-plane.  The signs of $B_{2j}(1/2)$ are $(-1)^j$, including $j=0$.
Every summand in \eqref{tp:eq:h} therefore has sign $(-1)^j$; hence
$h_{d,j}\ne0$ and the listed poles are actual.

At $s=-m$ the terms subtracted in \eqref{tp:eq:ordinary} contribute
$h_{d,j}\zeta(d+2j-m,1/2)=0$.  Indeed their Hurwitz arguments are
nonpositive even integers, and $\zeta(-2q,1/2)=0$ for $q\geq0$.
The next summand is $O(x_n^{-2})$, or decays faster when no subtraction
is needed.  This proves both the ordinary convergence and the equality.
\end{proof}

Thus the word ``moment'' below always refers to the explicit convergent
expression \eqref{tp:eq:ordinary}.  The meromorphic family is a proof
and differentiation tool; no assignment of a value to an unspecified
divergent sum is required.

\subsection{Finite-cutoff coordinates}

For a finite-cutoff asymptotic expansion, let $\CT_N$ denote the coefficient
of $N^0(\log N)^0$.  Put
\begin{equation}\label{tp:eq:TC}
 \begin{split}
 T_{e,q}&=\CT_N\sum_{n=1}^N n^q A_n^e,\qquad
 b_e(k)=[u^k]\mathcal A(u)^e,\\
 \mathcal A(u)&=\sum_{a\geq1}B_{a-1}u^a,\\
 C_{d,k}&=b_d(k)+\sum_{j=1}^d\binom dj T_{d-j,k-2j}.
 \end{split}
\end{equation}
Here $b_0(k)=\delta_{k,0}$, and a coefficient with negative index is
zero.  The series $\mathcal A$ is formal.  These constant coefficients
exist: repeated finite summation by parts reduces their existence to
ordinary power sums, harmonic numbers, and convergent multiple zeta
values.  A constructive reduction is given below.

The exact recurrence
$A_{n-1}=A_n+n^{-2}$ shows that $C_{d,k}$ is the constant coefficient
of
\[
 \sum_{n=1}^N n^k(A_{n-1}^d-A_n^d)+N^kA_N^d.
\]
Define the centered Faulhaber polynomial
\begin{equation}\label{tp:eq:Faulhaber}
 P_m(N)=\sum_{n=0}^{N-1}(n+1/2)^m
 =\frac{B_{m+1}(N+1/2)-B_{m+1}(1/2)}{m+1}.
\end{equation}
Finite summation by parts gives
\begin{equation}\label{tp:eq:finiteSBP}
 \sum_{n=0}^{N-1}x_n^m A_n^d
 =P_m(N)A_N^d+
   \sum_{n=1}^N P_m(n)(A_{n-1}^d-A_n^d).
\end{equation}
Every subtracted even-power sum in \eqref{tp:eq:ordinary} is an odd
polynomial in $N$, so it has zero constant coefficient.  Consequently,
if $P_m(N)=\sum_{k\geq1}p_{m,k}N^k$, then
\begin{equation}\label{tp:eq:momentC}
 M_{d,m}=\sum_{k\geq1}p_{m,k}C_{d,k}.
\end{equation}

The parity in \eqref{tp:eq:Faulhaber} is useful.  For even $m$ the
polynomial is odd; for odd $m$ it is even with zero constant term.
Formally define
\begin{equation}\label{tp:eq:GH}
 G_d(t)=\sum_{\substack{m\geq0\\m\equiv d\ (2)}}M_{d,m}\frac{t^m}{m!},
 \qquad H_d(t)=2\sinh(t/2)G_d(t).
\end{equation}
The finite polynomial identity \eqref{tp:eq:momentC} is equivalent to
\begin{equation}\label{tp:eq:HC}
 H_d(t)=\sum_{\substack{k\geq1\\k\equiv d+1\ (2)}}C_{d,k}\frac{t^k}{k!}.
\end{equation}
For example, summing \eqref{tp:eq:Faulhaber} over even $m$ gives
$\sinh(Nt)/(2\sinh(t/2))$; summing over odd $m$ gives
$(\cosh(Nt)-1)/(2\sinh(t/2))$.  These two identities prove
\eqref{tp:eq:HC} coefficient by coefficient without assuming convergence
of a moment generating series in advance.

\subsection{The universal boundary cancellation}

\begin{lemma}[A finite Bernoulli difference identity]
\label{tp:lem:EM}
For every integer $d\geq1$ and indeterminates $a,h$,
\begin{equation}\label{tp:eq:EM}
 \begin{split}
 (a+h)^d-a^d={}&\frac{dh}{2}\bigl((a+h)^{d-1}+a^{d-1}\bigr)\\
 &-\sum_{j=1}^{\lfloor(d-1)/2\rfloor}
      \binom d{2j}B_{2j}h^{2j}
      \bigl((a+h)^{d-2j}-a^{d-2j}\bigr).
 \end{split}
\end{equation}
Moreover, whenever $k\equiv d+1\pmod2$,
\begin{equation}\label{tp:eq:boundary}
 b_d(k)+\frac d2 b_{d-1}(k-2)
 +\sum_{j=1}^{\lfloor d/2\rfloor}
      \binom d{2j}B_{2j}b_{d-2j}(k-4j)=0.
\end{equation}
\end{lemma}

\begin{proof}
The first identity is the terminating Euler--Maclaurin identity for a
polynomial on one interval.  It also follows by comparing coefficients
in the Bernoulli generating function, so its use here is purely finite
algebra.  A term with $d-2j=0$ would be a difference of constants and is
omitted.

For the second identity write
\[
 E(u)=\sum_{j\geq0}B_{2j}u^{2j},\qquad
 \mathcal A(u)=uE(u)-u^2/2.
\]
The expression on the left of \eqref{tp:eq:boundary} is the coefficient
of $u^k$ in
\begin{equation}\label{tp:eq:BernoulliMechanism}
 \begin{split}
 &\mathcal A(u)^d+\frac d2u^2\mathcal A(u)^{d-1}
 +\sum_{j\geq1}\binom d{2j}B_{2j}u^{4j}\mathcal A(u)^{d-2j}\\
 &\hspace{32mm}=u^{2d}B_d\left(\frac12+\frac{E(u)}u\right).
 \end{split}
\end{equation}
The Bernoulli reflection law $B_d(1-z)=(-1)^dB_d(z)$ implies that
the right side has parity $d$ as a formal series in $u$.  Its
coefficient at the opposite parity is zero.  Although the argument of
$B_d$ contains $u^{-1}$, the displayed product has no negative powers:
the polynomial has degree $d$.  This establishes the cancellation.
\end{proof}

Identity \eqref{tp:eq:BernoulliMechanism} supplies a single reason for
all the boundary terms to cancel.  The separate quartic cancellation
is one low-degree consequence.

\begin{theorem}[All-power triangular recurrence]
\label{tp:thm:recurrence}
Set $G_0=0$.  For $d\geq1$ and $k\geq2$ with
$k\equiv d+1\pmod2$,
\begin{equation}\label{tp:eq:recurrence}
 \boxed{\begin{aligned}
 C_{d,k}={}&d(k-2)![t^{k-2}]\bigl(\cosh(t/2)G_{d-1}(t)\bigr)
       -\frac d2\delta_{k,2}\zeta(2)^{d-1}\\
 &-\sum_{j=1}^{\lfloor(d-1)/2\rfloor}
       \binom d{2j}B_{2j}C_{d-2j,k-4j}.
 \end{aligned}}
\end{equation}
Every coordinate on the right has smaller tail power.
\end{theorem}

\begin{proof}
Apply \eqref{tp:eq:EM} with $a=A_n$ and $h=n^{-2}$, multiply by
$n^k$, and sum only up to a finite $N$.  For $m=k-2\geq0$, changing
the index in one of the two symmetric sums gives the exact relation
\begin{align*}
 &\sum_{n=1}^N n^m(A_{n-1}^{d-1}+A_n^{d-1})\\
 &\quad=\sum_{n=0}^{N-1}\bigl((n+1)^m+n^m\bigr)A_n^{d-1}
       -\delta_{m,0}\zeta(2)^{d-1}+N^mA_N^{d-1}.
\end{align*}
The polynomial $(x+1/2)^m+(x-1/2)^m$ contains only powers of
parity $m\equiv d-1\pmod2$.  Taking its constant coefficient therefore
uses exactly the regular moments of power $d-1$.  Its contribution is
$2m![t^m](\cosh(t/2)G_{d-1})$.  For $d=1$, these polynomial sums have
zero constant coefficient, which is the convention $G_0=0$.

The differences of the lower powers in \eqref{tp:eq:EM} give
$C_{d-2j,k-4j}-b_{d-2j}(k-4j)$.  The remaining boundary terms are
precisely \eqref{tp:eq:boundary}; their sum is zero.  Taking constant
coefficients after these finite identities have been combined proves
\eqref{tp:eq:recurrence}.  In particular the exceptional $k=2$ term
comes from the missing $n=0$ endpoint and cannot be discarded.
\end{proof}

\subsection{Explicit anchors and a convergent polylogarithm generator}

Write $\boldsymbol a\vDash e$ for a composition of $e$.  For $e\geq1$
and $h\leq1$, excluding $(e,h)=(1,1)$, define the finite MZV combination
\begin{equation}\label{tp:eq:anchor}
 \mathcal Z_{e,h}=
 \sum_{(a_1,\ldots,a_\ell)\vDash e}
 \frac{e!}{a_1!\cdots a_\ell!}
 \zeta(2a_1,\ldots,2a_{\ell-1},2a_\ell-h).
\end{equation}
All the displayed multiple zeta values converge.  Ordering the indices
by their minimum, as in Hoffman's tail-product method \cite{HoffmanTails},
gives
\begin{equation}\label{tp:eq:anchorseries}
 C_{e,h}=\mathcal Z_{e,h}
 =\sum_{n\geq1}n^h(A_{n-1}^e-A_n^e),\qquad h\leq1.
\end{equation}
The exceptional pair is divergent and is never needed.  For $h=0$,
telescoping yields $\mathcal Z_{e,0}=\zeta(2)^e$.  Formula
\eqref{tp:eq:anchor} is classical index ordering, not a new claim about
unweighted tail products.

Let $I^q$ denote an anchored $q$-fold primitive:
\begin{equation}\label{tp:eq:primitive}
 (I^qf)(t)=\frac1{(q-1)!}\int_0^t(t-u)^{q-1}f(u)\,du,
 \qquad q\geq1.
\end{equation}
For $d\geq1$ put $c_{d,j}=\binom d{2j}B_{2j}$ and let
\begin{equation}\label{tp:eq:anchorpoly}
 \begin{split}
 R_d(t)={}&\mathbf1_{2\mid d}\,\mathcal Z_{d,1}t
       -\mathbf1_{2\nmid d}\,\frac d4\zeta(2)^{d-1}t^2\\
 &-\sum_{j=1}^{\lfloor(d-1)/2\rfloor}c_{d,j}
       \sum_{\substack{2\leq k\leq4j\\k\equiv d+1\ (2)}}
         \mathcal Z_{d-2j,k-4j}\frac{t^k}{k!}.
 \end{split}
\end{equation}
For $d=1$ the first term is absent; it is not an invocation of the
undefined anchor $\mathcal Z_{1,1}$.

\begin{theorem}[An anchored generator at every power]
\label{tp:thm:generator}
The series in \eqref{tp:eq:GH} converge for $|t|<2\pi$.  They are
determined recursively by $G_0=H_0=0$ and
\begin{equation}\label{tp:eq:generator}
 \boxed{\begin{aligned}
 H_d(t)={}&R_d(t)+dI^2\bigl(\cosh(t/2)G_{d-1}(t)\bigr)\\
 &-\sum_{j=1}^{\lfloor(d-1)/2\rfloor}c_{d,j}I^{4j}H_{d-2j}(t),
 \qquad G_d(t)=\frac{H_d(t)}{2\sinh(t/2)}.
 \end{aligned}}
\end{equation}
The apparent singularity at zero is removable.  At each fixed $d$,
all moments $M_{d,m}$ lie in one explicitly specified finite rational
span of convergent MZVs and $1$, independent of $m$.
One uniform, possibly redundant choice is the span of $1$ and all
positive-index MZVs of weight at most $2d-1$ with first index at least two.

Moreover $H_d$ has a finite expression in polynomials in $t$, multiple
zeta constants, and one-variable multiple polylogarithms at $e^{-t}$
of functional depth at most $d-1$.  This is an upper bound in a
specified representation, not a minimal-depth or period-independence claim.
\end{theorem}

\begin{proof}
Sum \eqref{tp:eq:recurrence} with the factors $t^k/k!$.  The positive
lower indices give $I^{4j}H_{d-2j}$.  The nonpositive lower indices
give exactly the finite polynomial in \eqref{tp:eq:anchorpoly}, and
the $k=1$ coefficient at even $d$ is \eqref{tp:eq:anchorseries}.  Thus
\eqref{tp:eq:generator} holds first as a formal power-series identity.

Inductively the numerator is holomorphic for $|t|<2\pi$ and has
parity $d+1$.  It is $O(t)$ at even $d$ and $O(t^2)$ at odd $d$.
The denominator has only a simple zero at zero inside this disk.
Division therefore preserves holomorphy and gives parity $d$.
This proves convergence of the formal moment series and the asserted
value of its coefficients.  At fixed $d$ only finitely many anchors
from powers at most $d$ have been introduced; their weights are at
most $2d-1$.  Taylor coefficients of all remaining elementary factors
and integrations are rational.  This proves the finite-span assertion.

For the functional representation start with real $t>0$ and put
$z=e^{-t}$. The assembled expression will then denote the analytic
continuation supplied by the anchored generator on $|t|<2\pi$;
separated terms need not be holomorphic at zero. The factor that
enters the new integration step is
\[
 \cosh(t/2)G_{d-1}(t)
 =\left(\frac12+\frac z{1-z}\right)H_{d-1}(t).
\]
Use the convention
\[
 L_{a_1,\ldots,a_\ell}(t)
 =\sum_{n_1>\cdots>n_\ell\geq1}
       \frac{e^{-tn_1}}{n_1^{a_1}\cdots n_\ell^{a_\ell}}.
\]
The empty word represents $1$. We maintain the stronger induction
invariant that each nonempty polylogarithm in $H_d$ has first index at
least two. Ordinary integrations raise that first index. For the new
operation, an explicit double primitive of
$t^m zL_{\boldsymbol a}(t)/(1-z)$ is
\[
 P_{m,\boldsymbol a}(t)=\sum_{j=0}^m(j+1)\frac{m!}{(m-j)!}
                t^{m-j}L_{j+2,\boldsymbol a}(t).
\]
Indeed $L_{1,\boldsymbol a}'=-zL_{\boldsymbol a}/(1-z)$ and
$L_{a+1,\boldsymbol b}'=-L_{a,\boldsymbol b}$. This raises depth by
at most one and preserves the leading-index invariant. Its value at
zero is a convergent MZV, and its first derivative is
\[
 P_{m,\boldsymbol a}'(t)=-\sum_{j=0}^m\frac{m!}{(m-j)!}
                t^{m-j}L_{j+1,\boldsymbol a}(t).
\]
For $m>0$ the endpoint is $-m!\zeta(m+1,\boldsymbol a)$. For $m=0$
the only potential divergence has exactly one leading index one.

For an admissible word $\boldsymbol a=(a_1,\ldots,a_\ell)$ define
\begin{align*}
 K(\boldsymbol a)={}&-\sum_{i=1}^{\ell}
 \zeta(a_1,\ldots,a_i,1,a_{i+1},\ldots,a_\ell)\\
 &-\sum_{i=1}^{\ell}
 \zeta(a_1,\ldots,a_i+1,\ldots,a_\ell),
 \qquad K(\varnothing)=0.
\end{align*}
All these values converge, and
\[
 L_{1,\boldsymbol a}(t)
       =\zeta(\boldsymbol a)L_1(t)+K(\boldsymbol a)+o(1).
\]
To prove this, subtract $\zeta(\boldsymbol a)L_1(t)$ termwise.
The remaining series is
$\sum_{n\ge1}e^{-tn}(H_{n-1}(\boldsymbol a)-\zeta(\boldsymbol a))/n$,
where $H_{n-1}(\boldsymbol a)$ is the finite multiple harmonic sum.
It is absolutely convergent at $t=0$ because $a_1\ge2$. Ordering its
extra index among the zeta indices gives exactly the two sums defining
$K$. Dominated convergence proves the displayed endpoint formula.

In the complete double primitive the coefficient of the potentially
divergent $L_1(t)$ is $-H_{d-1}(0)=0$. Every remaining endpoint is
therefore an explicit convergent MZV. Subtracting the value and first
derivative at zero anchors the primitive. The other integrations in
\eqref{tp:eq:generator} do not increase depth. This proves the bound
$d-1$ and fixes every integration constant without assigning values
to individual divergent multiple zeta series.
\end{proof}

For completeness, if $F_j'=-F_{j-1}$, an unanchored $q$-fold primitive
of $t^mF_0(t)$ is
\begin{equation}\label{tp:eq:finiteprimitive}
 (-1)^q\sum_{j=0}^m\binom{q+j-1}{j}
       \frac{m!}{(m-j)!}t^{m-j}F_{q+j}(t).
\end{equation}
Subtracting its Taylor polynomial of degree $q-1$ at zero supplies
the anchor, with grouped limits when necessary.  Thus
\eqref{tp:eq:generator} is also a finite algorithm for antiderivatives
and multiple-polylogarithm formulas, rather than a recurrence with
unspecified constants.

\subsection{A finite-cutoff implementation independent of the new recurrence}

The supplied verification program does not use \eqref{tp:eq:generator}
to construct the comparison moments.  If
$F_q(N)=\sum_{n=1}^N n^q=\sum_a f_aN^a$, ordinary finite summation
by parts instead gives
\begin{equation}\label{tp:eq:independentT}
 T_{e,q}=\sum_a f_ab_e(a)+
 \sum_{j=1}^e\binom ej
 \left(\sum_a f_aT_{e-j,a-2j}-T_{e-j,q-2j}\right),\qquad q\geq0.
\end{equation}
For $q<0$, products of strictly greater tail indices give
\[
 T_{e,q}=\sum_{\boldsymbol a\vDash e}
       \frac{e!}{a_1!\cdots a_\ell!}
       \zeta(2a_1,\ldots,2a_\ell,-q),\qquad e\geq1.
\]
The remaining base cases are $T_{0,q}=0$ for $q\geq0$,
$T_{0,-1}=\gamma$, and $T_{0,-p}=\zeta(p)$ for $p\geq2$.
The power $e$ decreases at every recursive step.  This proves existence
of the constant coefficients used above and supplies a terminating
independent exact verification route.

\subsection{The same cancellation for every even Hurwitz order}
\label{et:section}

Fix $q\geq1$ and put $p=2q$. In this subsection the symbols
$M,C,G,H$ refer to the fixed tail
$T_n=\zeta(p,n+1)=\psi^{(p-1)}(n+1)/(p-1)!$, with
$x_n=n+1/2$. Thus $T_{n-1}=T_n+n^{-p}$. Define
\[
 \eta_0=\frac1{p-1},\qquad
 \eta_j=\frac{(p)_{2j-1}}{(2j)!}B_{2j}(1/2)\quad(j\geq1),
 \qquad h_{d,j}=\sum_{j_1+\cdots+j_d=j}\prod_{\nu=1}^d\eta_{j_\nu}.
\]
The centered Hurwitz asymptotics give
$T_n^d\sim\sum_{j\geq0}h_{d,j}x_n^{-(p-1)d-2j}$.
For $m\geq0$, $m\equiv d\pmod2$, the regular continued moment is
the ordinary absolutely convergent sum
\begin{equation}\label{et:ordinary}
 M_{d,m}=\sum_{n\geq0}\left\{x_n^mT_n^d-
 \sum_{0\leq j\leq(m-(p-1)d)/2}
              h_{d,j}x_n^{m-(p-1)d-2j}\right\}.
\end{equation}
Indeed the next term is $O(x_n^{-2})$, while every removed term
has continued sum $\zeta(-2r,1/2)=0$. The meromorphic family
$\sum T_n^dx_n^{-s}$ has its poles at $1-(p-1)d-2j$;
the same-sign convolution of the $\eta_j$ shows that their residues
$h_{d,j}$ are nonzero.

For every integer $k$ put
\[
 C_{d,k}=\CT_N\left\{\sum_{n=1}^N n^k(T_{n-1}^d-T_n^d)
                                  +N^kT_N^d\right\},\qquad
 G_d(t)=\sum_{m\equiv d\ (2)}M_{d,m}\frac{t^m}{m!}.
\]
Subtracting finitely many asymptotic terms proves existence of each
cutoff constant. The latter sum has $m\geq0$.
Set $H_d=2\sinh(t/2)G_d$ and
$G_0=H_0=0$. The centered Faulhaber identity used above gives
\begin{equation}\label{et:HC}
 H_d(t)=\sum_{\substack{k\geq1\\k\equiv d+1\ (2)}}C_{d,k}\frac{t^k}{k!}.
\end{equation}
For $e\geq1$ and $h<p$, except $(e,h)=(1,p-1)$, the needed anchors are
\begin{equation}\label{et:anchors}
 \mathcal Z^{(q)}_{e,h}=
 \sum_{(a_1,\ldots,a_\ell)\vDash e}
 \frac{e!}{a_1!\cdots a_\ell!}
 \zeta(pa_1,\ldots,pa_{\ell-1},pa_\ell-h)=C_{e,h}.
\end{equation}
Here $\vDash$ denotes a composition. Ordering the $e$ tail indices
by their minimum proves this formula: the summand is
$n^h(T_{n-1}^e-T_n^e)=O(n^{h-(p-1)e-1})$ and the boundary term
tends to zero. The excluded pair never has the required parity.

Write $c_{d,j}=\binom d{2j}B_{2j}$ and define the finite polynomial
\begin{align}\label{et:R}
 R_d(t)={}&\sum_{\substack{1\leq k<p\\k\equiv d+1\ (2)}}
       \mathcal Z^{(q)}_{d,k}\frac{t^k}{k!}
 -\mathbf1_{2\nmid d}\frac d2\zeta(p)^{d-1}\frac{t^p}{p!}\notag\\
 &-\sum_{j=1}^{\lfloor(d-1)/2\rfloor}c_{d,j}
   \sum_{\substack{p\leq k\leq2pj\\k\equiv d+1\ (2)}}
       \mathcal Z^{(q)}_{d-2j,k-2pj}\frac{t^k}{k!}.
\end{align}

\begin{theorem}[Every power of every even-order tail]\label{et:main}
For $k\geq p$ with $k\equiv d+1\pmod2$,
\begin{equation}\label{et:recurrence}
 \begin{split}
 C_{d,k}={}&d(k-p)![t^{k-p}]\bigl(\cosh(t/2)G_{d-1}(t)\bigr)
       -\frac d2\delta_{k,p}\zeta(p)^{d-1}\\
 &-\sum_{j=1}^{\lfloor(d-1)/2\rfloor}c_{d,j}C_{d-2j,k-2pj}.
 \end{split}
\end{equation}
With the anchored primitives of \eqref{tp:eq:primitive}, its generator is
\begin{equation}\label{et:generator}
 H_d=R_d+dI^p\bigl(\cosh(t/2)G_{d-1}\bigr)
        -\sum_{j=1}^{\lfloor(d-1)/2\rfloor}c_{d,j}I^{2pj}H_{d-2j},
 \qquad G_d=\frac{H_d}{2\sinh(t/2)}.
\end{equation}
Both functions are holomorphic for $|t|<2\pi$. At fixed $(q,d)$,
all the moments belong to a fixed finite rational span of $1$ and
convergent MZVs of weights at most $2qd-1$. The generator $H_d$ admits
a finite multiple-polylogarithm representation of functional depth at
most $d-1$, with all constants fixed by the anchored primitives.
\end{theorem}

\begin{proof}
The uncentered formal asymptotic series is
\[
 \mathcal A_q(u)=u^{p-1}E_q(u)-\tfrac12u^p,
 \qquad E_q(u)=\frac1{p-1}
    +\sum_{j\geq1}\frac{(p)_{2j-1}B_{2j}}{(2j)!}u^{2j}.
\]
Its boundary combination satisfies the exact formal identity
\begin{equation}\label{et:boundary}
 \mathcal A_q^d+\frac d2u^p\mathcal A_q^{d-1}
 +\sum_{j=1}^{\lfloor d/2\rfloor}c_{d,j}u^{2pj}\mathcal A_q^{d-2j}
 =u^{pd}B_d\left(\frac12+\frac{E_q(u)}u\right).
\end{equation}
Since $E_q$ is even, Bernoulli reflection makes the right side have
parity $d$. Thus the required opposite-parity coefficients vanish.
Apply the finite polynomial identity \eqref{tp:eq:EM} with
$a=T_n$, $h=n^{-p}$ before taking constant coefficients. For
$r=k-p\geq0$, the symmetric term has the exact cutoff shift
\[
 \sum_{n=1}^N n^r(T_{n-1}^{d-1}+T_n^{d-1})
 =\sum_{n=0}^{N-1}\bigl((n+1)^r+n^r\bigr)T_n^{d-1}
   +N^rT_N^{d-1}-\delta_{r,0}\zeta(p)^{d-1}.
\]
The centered bracket has regular parity and gives
$2r![t^r](\cosh(t/2)G_{d-1})$. Equation~\eqref{et:boundary}
cancels the remaining boundaries, proving \eqref{et:recurrence}.
For $d=1$ the polynomial sums have zero constant coefficient, in
agreement with $G_0=0$.

Summing the recurrence gives \eqref{et:generator}: the low positive
indices give the first term of $R_d$, and the nonpositive lower
indices give its last term. Induction gives holomorphy on the stated
disk; $H_d$ has parity $d+1$ and vanishes at zero, so the division is
removable. Each low anchor has weight $pe-h\leq pd-1$; each last-term
anchor has weight $p(d-2j)-(k-2pj)=pd-k\leq pd-p$.
All other Taylor operations are rational. Finally
$\cosh(t/2)G_{d-1}=(1/2+z/(1-z))H_{d-1}$, $z=e^{-t}$;
the same primitive construction as in Theorem~\ref{tp:thm:generator}
raises functional depth by at most one. Increasing the number of
anchored integrations from $2$ to $p$ adds no depth.
\end{proof}

This extends the inspected two-tail and trigamma-power results to every
repeated even order. It does not yet address unequal orders among
three or more factors. No priority claim beyond the inspected source
corpus is intended.

\begin{corollary}[An ordinary cubic identity at every even order]
\label{et:cubic}
For every $q\geq1$, with $T_n=\zeta(2q,n+1)$,
\begin{equation}\label{et:cubic-formula}
 \boxed{\sum_{n\geq0}\left\{
       \bigl((n+1)^{2q}-n^{2q}\bigr)T_n^3-3T_n^2\right\}
       =-\frac{3\zeta(2q)^2+\zeta(4q)}2.}
\end{equation}
In particular, for $x_n=n+1/2$ and $T_n=\psi^{(3)}(n+1)/6$,
\begin{equation}\label{et:psi3}
 \boxed{\sum_{n\geq0}
       \bigl((4x_n^3+x_n)T_n^3-3T_n^2\bigr)
                      =-\frac94\zeta(8).}
\end{equation}
Both identities are absolutely convergent without subtraction.
\end{corollary}

\begin{proof}
Put $p=2q$. The cubic case of \eqref{tp:eq:EM}, followed by finite
summation by parts, gives the exact cutoff identity
\begin{align*}
 &\sum_{n=0}^{N-1}\left\{
       \bigl((n+1)^p-n^p\bigr)T_n^3-3T_n^2\right\}\\
 &\qquad=N^pT_N^3+\frac32T_N^2
       -\frac32\zeta(p)^2-\frac12H_N^{(2p)}.
\end{align*}
Here the cubic difference is
$n^p(T_{n-1}^3-T_n^3)
 =\tfrac32(T_{n-1}^2+T_n^2)-\tfrac12n^{-2p}$.
Since $T_n=O((n+1)^{1-p})$, the two boundary terms tend to zero
and each series on the left converges absolutely for $p\geq2$.
Letting $N\to\infty$ proves \eqref{et:cubic-formula}.
For $p=4$, use $(n+1)^4-n^4=4x_n^3+x_n$ and
$\zeta(4)^2=7\zeta(8)/6$ to obtain \eqref{et:psi3}.
\end{proof}

\begin{corollary}[Complex orders and a critical cubic value]
\label{et:complex-cubic}
Put $T_n(p)=\zeta(p,n+1)$, using the real logarithm of positive
integers in all powers and setting $0^p=0$ in the domain below.
For every complex $p$ with $\Re p>3/2$,
\begin{equation}\label{et:complex-cubic-formula}
 \sum_{n\geq0}\left\{
    \bigl((n+1)^p-n^p\bigr)T_n(p)^3-3T_n(p)^2\right\}
       =-\frac{3\zeta(p)^2+\zeta(2p)}2.
\end{equation}
At the real critical order $p=3/2$, the combined series also converges
absolutely, but its value has an additional boundary term:
\begin{equation}\label{et:critical-cubic}
 \boxed{\begin{aligned}
 &\sum_{n\geq0}\left\{
    \bigl((n+1)^{3/2}-n^{3/2}\bigr)T_n(3/2)^3-3T_n(3/2)^2\right\}
 \\[-1mm]
 &\qquad=8-\frac{3\zeta(3/2)^2+\zeta(3)}2.
 \end{aligned}}
\end{equation}
\end{corollary}

\begin{proof}
The preceding cutoff identity remains exact for $\Re p>1$.
For fixed $p$ in this half-plane,
$T_N(p)=N^{1-p}/(p-1)+O(N^{-\Re p})$.
If $\Re p>3/2$, each series on its left converges absolutely and
$N^pT_N(p)^3=O(N^{3-2\Re p})\to0$, proving
\eqref{et:complex-cubic-formula}.

At $p=3/2$, put $x=n+1/2$. The centered expansions are
\[
 T_n(3/2)=2x^{-1/2}-\frac1{16}x^{-5/2}+O(x^{-9/2}),\qquad
 (n+1)^{3/2}-n^{3/2}
       =\frac32x^{1/2}-\frac1{64}x^{-3/2}+O(x^{-7/2}).
\]
Consequently the combined summand in \eqref{et:critical-cubic} is
$-\tfrac12x^{-3}+O(x^{-5})$, proving absolute convergence.
In the exact cutoff formula,
$N^{3/2}T_N(3/2)^3\to8$, $T_N(3/2)^2\to0$, and
$H_N^{(3)}\to\zeta(3)$. These limits give the stated critical value.
\end{proof}


\section{Fifth and sixth powers: finite coordinate spaces and explicit identities}
\label{tp:sec:examples}

We now make the all-power theorem concrete.  Retain the classical
constants
\begin{align}
 Q_5&=4\zeta(2)\zeta(3)-10\zeta(5),\label{tp:eq:Q5}\\
 Q_7&=51\zeta(3)\zeta(4)+10\zeta(2)\zeta(5)
                       -\frac{301}{4}\zeta(7).\label{tp:eq:Q7}
\end{align}
The latter is $\sum_{n\geq0}A_n^4$, by Hoffman's quartic evaluation
with the $n=0$ term included \cite{HoffmanTails,RepoQuartic}.

\subsection{The lower-power inputs}

The universal recurrence starts with
\begin{align}
 H_1(t)&=-\frac{t^2}{4},\label{tp:eq:H1}\\
 H_2(t)&=3\zeta(3)t+2I^2(\cosh(t/2)G_1(t)),\label{tp:eq:H2}\\
 H_3(t)&=3I^2(\cosh(t/2)G_2(t))
      -\frac{17}{8}\zeta(4)t^2-\frac{\zeta(2)}{48}t^4
      +\frac{t^6}{2880},\label{tp:eq:H3}\\
 H_4(t)&=Q_7t+\frac{Q_5t^3}{6}
      +4I^2(\cosh(t/2)G_3(t))-I^4H_2(t).\label{tp:eq:H4}
\end{align}
These reproduce the existing quadratic, cubic, and quartic generators.
For example $\mathcal Z_{2,-1}=2\zeta(2,3)+\zeta(5)=-Q_5$ and
$\mathcal Z_{e,0}=\zeta(2)^e$ account for their exceptional terms.

\subsection{The quintic generator and its eight coordinates}

Define the one new weight-eight combination
\begin{equation}\label{tp:eq:K8}
 K_8=\frac56\mathcal Z_{3,-2}+\frac{289}{32}\zeta(8),
 \quad
 \mathcal Z_{3,-2}=
 6\zeta(2,2,4)+3\zeta(2,6)+3\zeta(4,4)+\zeta(8).
\end{equation}
This is a specified convergent MZV coordinate, not a claim of a new
transcendental constant or of its irreducibility.

\begin{theorem}[Every regular quintic moment]
\label{tp:thm:five}
For $|t|<2\pi$,
\begin{equation}\label{tp:eq:H5}
 \boxed{\begin{aligned}
 H_5(t)={}&5I^2(\cosh(t/2)G_4(t))
            -\frac53I^4H_3(t)-K_8t^2
            -\frac{19}{64}\zeta(6)t^4\\
          &+\frac{\zeta(4)}{4320}t^6
            +\frac{\zeta(2)}{241920}t^8
            -\frac{t^{10}}{43545600}.
 \end{aligned}}
\end{equation}
All the moments $M_{5,2r+1}$ lie in
\begin{equation}\label{tp:eq:space5}
 \operatorname{span}_{\mathbb Q}
 \{K_8,\zeta(6),Q_7,Q_5,\zeta(2),\zeta(3),\zeta(4),1\}.
\end{equation}
The coefficient matrix of the first eight moments in this ordered
list has determinant
\begin{equation}\label{tp:eq:rank5}
 \frac{672072587373661}{748313054183581286400}\ne0.
\end{equation}
Thus eight is the exact rank of these specified rational coefficient
rows.  No arithmetic independence of the eight evaluated constants
is asserted.
\end{theorem}

\begin{proof}
At $d=5$, the two Bernoulli coefficients in \eqref{tp:eq:generator}
are $5/3$ and $-1/6$.  Insert the anchors
$\mathcal Z_{3,-2}$, $\mathcal Z_{3,0}=\zeta(2)^3$, and
$\mathcal Z_{1,-2j}=\zeta(2j+2)$.  Euler's even-zeta evaluations
give $\zeta(2)^3=35\zeta(6)/8$ and
$\zeta(2)^4=175\zeta(8)/24$.  Finally insert
$I^8H_1=-t^{10}/7257600$.  Collecting these finite terms gives
\eqref{tp:eq:H5}.  The span follows from the lower-power generators.
Exact rational coefficient extraction gives \eqref{tp:eq:rank5};
the complete matrix and an executable determinant calculation accompany
the article.
\end{proof}

The first three formulas, useful independently of the generator, are
\begin{align}
 M_{5,1}={}&-K_8+\frac52Q_7,\label{tp:eq:M51}\\
 M_{5,3}={}&\frac14K_8+\frac5{12}Q_5-\frac5{12}Q_7
           +\frac{15}{2}\zeta(3)-\frac{85}{24}\zeta(4)
           -\frac{57}{32}\zeta(6),\label{tp:eq:M53}\\
 M_{5,5}={}&-\frac7{48}K_8-\frac5{72}Q_5+\frac{47}{288}Q_7
           -\frac1{12}\zeta(2)-\frac52\zeta(3)
           +\frac{157}{144}\zeta(4)\notag\\
          &+\frac{95}{64}\zeta(6)-\frac1{12}.
          \label{tp:eq:M55}
\end{align}
The first two are unsubtracted sums.  The third is
$\sum_{n\geq0}(x_n^5A_n^5-1)$.

\begin{corollary}[Two ordinary quintic evaluations]
\label{tp:cor:explicit}
With $x=x_n$ and $A=A_n$, both series below converge absolutely:
\begin{align}
 &\sum_{n\geq0}\bigl((23x+232x^3+240x^5)A^5-240\bigr)\notag\\
 &\qquad=80Q_5-20\zeta(2)+1140\zeta(3)-560\zeta(4)
                         -57\zeta(6)-20,\label{tp:eq:short5}\\[3pt]
 &\boxed{\begin{aligned}
 &\sum_{n\geq0}\left[
 (195x+4056x^3+11568x^5+5376x^7)A^5
                  -5376x^2-9328\right]\\
 &\hspace{12mm}=-748\zeta(2)+7884\zeta(3)-5380\zeta(4)
                  +171\zeta(6)-748.
 \end{aligned}}\label{tp:eq:main5}
\end{align}
\end{corollary}

\begin{proof}
The first identity is $23M_{5,1}+232M_{5,3}+240M_{5,5}$.
Its $K_8$ and $Q_7$ coefficients cancel.  The second is
$195M_{5,1}+4056M_{5,3}+11568M_{5,5}+5376M_{5,7}$.
The three coefficients of $K_8,Q_7,Q_5$ cancel.  Its subtraction
polynomial is $5376(x^2-5/12)+11568=5376x^2+9328$.

There is a particularly small independent MZV check of the second
formula.  The finite-cutoff recursion \eqref{tp:eq:independentT}
reduces its left side to
\begin{align*}
 &-748\zeta(2)+5256\zeta(2,1)-3840\zeta(2,2)
   +240\zeta(2,2,2)\\
 &\quad+120\bigl(\zeta(2,4)+\zeta(4,2)\bigr)
   +2628\zeta(3)-2500\zeta(4)+36\zeta(6)-748.
\end{align*}
Use $\zeta(2,1)=\zeta(3)$,
$\zeta(2,2)=3\zeta(4)/4$,
$\zeta(2,2,2)=3\zeta(6)/16$, and
$\zeta(2,4)+\zeta(4,2)=3\zeta(6)/4$.
These are Euler, all-twos, and elementary stuffle identities.
The displayed result follows.  Absolute convergence follows directly
from \eqref{tp:eq:centeredasy}.
\end{proof}

The values of \eqref{tp:eq:short5} and \eqref{tp:eq:main5} begin,
respectively,
\[
 456.5518108340607851715385608386\ldots,
 \qquad1849.672608741185304145251799308\ldots.
\]
These decimals are diagnostic illustrations, not inputs to the proof.

\subsection{The sextic generator and eleven coordinates}

Define
\begin{equation}\label{tp:eq:new6}
 R_{11}=\mathcal Z_{6,1},\qquad
 K_9=\mathcal Z_{2,-5}-5\mathcal Z_{4,-1},\qquad
 J_7=\mathcal Z_{2,-3}=2\zeta(2,5)+\zeta(7).
\end{equation}
Formula \eqref{tp:eq:anchor} specifies every term of these finite
convergent MZV combinations.

\begin{theorem}[Every regular sextic moment]
\label{tp:thm:six}
For $|t|<2\pi$,
\begin{equation}\label{tp:eq:H6}
 \boxed{\begin{aligned}
 H_6(t)={}&R_{11}t+6I^2(\cosh(t/2)G_5(t))
          -\frac52I^4H_4(t)+\frac12I^8H_2(t)\\
         &+\frac{K_9}{12}t^3+\frac{J_7}{240}t^5
          -\frac{Q_5}{10080}t^7.
 \end{aligned}}
\end{equation}
All $M_{6,2r}$ belong to the rational span of the ordered list
\begin{equation}\label{tp:eq:space6}
 (R_{11},K_9,J_7,K_8,\zeta(6),Q_7,Q_5,
                       \zeta(2),\zeta(3),\zeta(4),1).
\end{equation}
The first eleven coefficient rows have determinant
\begin{equation}\label{tp:eq:rank6}
 -\frac{557325343494359097538722307}
         {135927478735611867783198720000000}\ne0.
\end{equation}
This is again a formal coefficient-rank statement.
\end{theorem}

\begin{proof}
The nonzero coefficients at $d=6$ in \eqref{tp:eq:generator} are
$5/2$ and $-1/2$.  The $k=1$ anchor is $R_{11}$; the remaining
negative-index anchors combine to $K_9,J_7$, and
$\mathcal Z_{2,-1}=-Q_5$.  This gives \eqref{tp:eq:H6} directly.
The span and determinant follow by rational coefficient extraction.
\end{proof}

For example,
\begin{align}
 M_{6,0}&=R_{11},\notag\\
 M_{6,2}&=-2K_8+\frac16K_9+5Q_7-\frac1{12}R_{11},\notag\\
 M_{6,4}&=\frac1{10}J_7+\frac25K_8-\frac1{12}K_9
       +\frac12Q_5-\frac54Q_7+\frac7{240}R_{11}\notag\\
       &\quad+9\zeta(3)-\frac{17}{4}\zeta(4)
                       -\frac{171}{80}\zeta(6).
                       \label{tp:eq:M6examples}
\end{align}
One further elimination gives an ordinary sixth-power identity:
\begin{equation}\label{tp:eq:main6}
 \begin{split}
 &\sum_{n\geq0}\left[
 (35+3360x^2+29792x^4+50176x^6+16128x^8)A^6
                    -16128x^2-42112\right]\\
 &\quad=9184Q_5-64Q_7-2928\zeta(2)+111216\zeta(3)
          -55368\zeta(4)-9576\zeta(6)-2928.
 \end{split}
\end{equation}
It is seven times the combination with coefficients
$(5,480,4256,7168,2304)$ of $M_{6,0},\ldots,M_{6,8}$.
Its value is $33200.31747342681779918605867754\ldots$.

\section{The resonant parity and all Laurent coefficients}
\label{tp:sec:resonance}

The regular moments do not determine spectral derivatives at the
interlaced poles.  The same explicit subtractions nevertheless give
the complete Laurent family, including all Stieltjes terms.  This
addresses the normalization part of the quartic report's Research
question R5 and applies to every power.

Use the convention
\begin{equation}\label{tp:eq:Stieltjes}
 \zeta(1+\varepsilon,a)=\frac1\varepsilon+
   \sum_{q\geq0}\frac{(-1)^q\gamma_q(a)}{q!}\varepsilon^q.
\end{equation}
For $d\geq1$, $r\geq0$, put $s_0=1-d-2r$ and
\[
 F_d(s_0+\varepsilon)=\frac{h_{d,r}}\varepsilon+
                 \sum_{q\geq0}\Lambda_{d,r,q}\frac{\varepsilon^q}{q!}.
\]

\begin{theorem}[Exact logarithmic sums at every resonance]
\label{tp:thm:Laurent}
For every $q\geq0$,
\begin{equation}\label{tp:eq:Laurent}
 \boxed{\begin{aligned}
 \Lambda_{d,r,q}={}&(-1)^q\sum_{n\geq0}(\log x_n)^q
 \left[x_n^{d+2r-1}A_n^d-
       \sum_{j=0}^{r}h_{d,j}x_n^{2r-1-2j}\right]\\
 &+\sum_{j=0}^{r-1}h_{d,j}
        \zeta^{(q)}(1+2j-2r,1/2)
       +(-1)^qh_{d,r}\gamma_q(1/2).
 \end{aligned}}
\end{equation}
The series is ordinary and absolutely convergent.  The Taylor series
of $F_d(s_0+\varepsilon)-h_{d,r}/\varepsilon$ has exact radius $2$.
\end{theorem}

\begin{proof}
The square bracket is $O(x_n^{-3})$. On a compact subset $K$ of
$\Re\varepsilon>-2$, its product with $x_n^{-\varepsilon}$ has
$q$th derivative bounded by
\[
 O_{K,q}\!\left(x_n^{-3-\min_{\varepsilon\in K}\Re\varepsilon}
                         (1+|\log x_n|)^q\right).
\]
This bound is summable, so the remainder and all fixed derivatives
converge normally there. Add
\[
 \sum_{j=0}^{r-1}h_{d,j}\zeta(1+2j-2r+\varepsilon,1/2)
 +h_{d,r}\bigl(\zeta(1+\varepsilon,1/2)-1/\varepsilon\bigr).
\]
This is the required pole-subtracted meromorphic family on a common
half-plane, hence by continuation near zero.  Differentiating there
gives \eqref{tp:eq:Laurent}, including the sign and normalization of
the Stieltjes term.  There are no poles in $|\varepsilon|<2$, while
$\varepsilon=-2$ is an actual pole with nonzero residue $h_{d,r+1}$.
This proves the radius assertion.
\end{proof}

At the first resonance the formula is especially short:
\begin{equation}\label{tp:eq:firstresonance}
 \Lambda_{d,0,q}=(-1)^q\left[
 \sum_{n\geq0}(\log x_n)^q
           \left(x_n^{d-1}A_n^d-x_n^{-1}\right)
          +\gamma_q(1/2)\right].
\end{equation}
All half-shift Stieltjes constants can be converted to the ordinary
constants $\gamma_q=\gamma_q(1)$ by
\begin{equation}\label{tp:eq:halfStieltjes}
 \gamma_q(1/2)=
 \frac{2(-1)^q(\log2)^{q+1}}{q+1}
 +2\sum_{j=0}^q\binom qj(-\log2)^{q-j}\gamma_j-\gamma_q.
\end{equation}
This follows by expanding
$\zeta(s,1/2)=(2^s-1)\zeta(s)$ at $s=1$.

The theorem gives a locally analytic, explicitly anchored generator in
the spectral displacement $\varepsilon$.  It does not identify a
finite arithmetic span as $r$ or the logarithmic order $q$ varies;
that is a further reduction problem.  In particular one must not
differentiate the discrete coefficient formulas for $M_{d,m}$ as
though they were identities at arbitrary complex $m$.
